% SoCC Revision: CPU Overhead Analysis
% Include this file in main.tex with: % SoCC Revision: CPU Overhead Analysis
% Include this file in main.tex with: % SoCC Revision: CPU Overhead Analysis
% Include this file in main.tex with: % SoCC Revision: CPU Overhead Analysis
% Include this file in main.tex with: \input{socc_cpu_overhead}

\subsection{CPU Overhead Analysis}\label{section:cpu_overhead}

Astrolabe's simulation runs entirely on CPU, separate from GPU-bound inference. We analyze CPU overhead on the 12-node A30 cluster, where each node has 32 CPU cores and runs 16 predictors.

\begin{figure}[tb]
    \centering
    \includegraphics[width=\linewidth]{figure/SoCC_revision/cpu_overhead.png}
    \caption{Predictor CPU utilization across QPS levels. Each predictor uses 21--32\% of a single core; 16 predictors consume $\sim$4 cores (12.5\% of node capacity).}
    \Description{CPU utilization measurements for Astrolabe predictor processes across QPS levels, showing per-predictor and aggregate CPU consumption on each node.}
    \label{fig:cpu_overhead}
\end{figure}

Figure~\ref{fig:cpu_overhead} shows per-predictor CPU usage ranges from 21\% to 32\% across QPS levels 20--32. With 16 predictors, aggregate consumption is $\sim$4 cores (12.5\% of 32-core capacity), leaving ample headroom for the vLLM frontend. Memory usage is $\sim$1.8~GB per predictor (29~GB total), well within the 128~GB available.Llumnix- and other heuristic schedulers require negligible CPU since they perform simple arithmetic rather than simulation.



\subsection{CPU Overhead Analysis}\label{section:cpu_overhead}

Astrolabe's simulation runs entirely on CPU, separate from GPU-bound inference. We analyze CPU overhead on the 12-node A30 cluster, where each node has 32 CPU cores and runs 16 predictors.

\begin{figure}[tb]
    \centering
    \includegraphics[width=\linewidth]{figure/SoCC_revision/cpu_overhead.png}
    \caption{Predictor CPU utilization across QPS levels. Each predictor uses 21--32\% of a single core; 16 predictors consume $\sim$4 cores (12.5\% of node capacity).}
    \Description{CPU utilization measurements for Astrolabe predictor processes across QPS levels, showing per-predictor and aggregate CPU consumption on each node.}
    \label{fig:cpu_overhead}
\end{figure}

Figure~\ref{fig:cpu_overhead} shows per-predictor CPU usage ranges from 21\% to 32\% across QPS levels 20--32. With 16 predictors, aggregate consumption is $\sim$4 cores (12.5\% of 32-core capacity), leaving ample headroom for the vLLM frontend. Memory usage is $\sim$1.8~GB per predictor (29~GB total), well within the 128~GB available.Llumnix- and other heuristic schedulers require negligible CPU since they perform simple arithmetic rather than simulation.



\subsection{CPU Overhead Analysis}\label{section:cpu_overhead}

Astrolabe's simulation runs entirely on CPU, separate from GPU-bound inference. We analyze CPU overhead on the 12-node A30 cluster, where each node has 32 CPU cores and runs 16 predictors.

\begin{figure}[tb]
    \centering
    \includegraphics[width=\linewidth]{figure/SoCC_revision/cpu_overhead.png}
    \caption{Predictor CPU utilization across QPS levels. Each predictor uses 21--32\% of a single core; 16 predictors consume $\sim$4 cores (12.5\% of node capacity).}
    \Description{CPU utilization measurements for Astrolabe predictor processes across QPS levels, showing per-predictor and aggregate CPU consumption on each node.}
    \label{fig:cpu_overhead}
\end{figure}

Figure~\ref{fig:cpu_overhead} shows per-predictor CPU usage ranges from 21\% to 32\% across QPS levels 20--32. With 16 predictors, aggregate consumption is $\sim$4 cores (12.5\% of 32-core capacity), leaving ample headroom for the vLLM frontend. Memory usage is $\sim$1.8~GB per predictor (29~GB total), well within the 128~GB available.Llumnix- and other heuristic schedulers require negligible CPU since they perform simple arithmetic rather than simulation.



\subsection{CPU Overhead Analysis}\label{section:cpu_overhead}

Astrolabe's simulation runs entirely on CPU, separate from GPU-bound inference. We analyze CPU overhead on the 12-node A30 cluster, where each node has 32 CPU cores and runs 16 predictors.

\begin{figure}[tb]
    \centering
    \includegraphics[width=\linewidth]{figure/SoCC_revision/cpu_overhead.png}
    \caption{Predictor CPU utilization across QPS levels. Each predictor uses 21--32\% of a single core; 16 predictors consume $\sim$4 cores (12.5\% of node capacity).}
    \Description{CPU utilization measurements for Astrolabe predictor processes across QPS levels, showing per-predictor and aggregate CPU consumption on each node.}
    \label{fig:cpu_overhead}
\end{figure}

Figure~\ref{fig:cpu_overhead} shows per-predictor CPU usage ranges from 21\% to 32\% across QPS levels 20--32. With 16 predictors, aggregate consumption is $\sim$4 cores (12.5\% of 32-core capacity), leaving ample headroom for the vLLM frontend. Memory usage is $\sim$1.8~GB per predictor (29~GB total), well within the 128~GB available.Llumnix- and other heuristic schedulers require negligible CPU since they perform simple arithmetic rather than simulation.

